\subsection{Exact reachability and a trial allowance after recovery}
\label{sec:elimination-reach}

Two costs limited the first batch implementation: repeated graph traversal
while planning, and an adaptive trial that could exhaust its entire allowance
before rebuilding a moderately large source presentation. The follow-up
preserves the greedy witness order and independent checker, but caches exact
reachability within each planning state and separates the continuation
allowance from presentation recovery. This changes resource use and some
finite-budget outcomes, not the certified Tietze operation.

\paragraph{Incremental cycle checks.}
At a planning state, a vertex has outgoing edges exactly when its definition
has already been selected. Let $R(v)$ be the set of vertices reachable from
$v$ by a nonempty directed path. Store these sets as dense integer masks and
store reverse adjacency for the selected edges. Both structures are fresh
for each planning state; only the earlier immutable word summaries persist
across states.

Suppose the next candidate defines $g$ using the dependency set $P$.
The generator $g$ has no previous outgoing edges. Form
\[
 D=\bigcup_{p\in P}\bigl(\{p\}\cup R(p)\bigr).
\]
Adding the candidate creates a directed cycle if and only if $g\in D$.
Reject it in that case. Otherwise set $R(g)=D$, add the reverse edges,
and propagate $D$ backwards through the old ancestors of $g$. An ancestor
already containing all of $D$ needs no further propagation.

For completeness, every new path must use one of the new edges from $g$.
Acyclicity prevents it from revisiting $g$, so its prefix is an old path to
$g$ and its suffix ends in $D$. Thus only old ancestors gain reachable
vertices, and all their new vertices belong to $D$. Conversely, every old
ancestor reaches all of $D$ after the update. If an ancestor $v$ already
contains $D$, each old ancestor of $v$ already contains it too, which
justifies pruning that propagation branch. Induction on accepted candidates
therefore maintains exact reachability. Since scoring, tie-breaking and
slot checks are unchanged, completed unconstrained runs select exactly the
same greedy batch as the earlier fresh depth-first tests.

Let $r$ be the number of mask positions, $k$ the accepted definitions and
$E$ their total dependency edges. If $\sum\delta$ counts dependency entries
examined in candidate queries, the query masks take $O(\sum\delta)$ mask
operations. During one accepted update each ancestor changes at most once;
reverse edges are then visited at most once for each changed endpoint.
Across the state, $O(kE)$ propagation visits are a conservative bound.
Allowing $O(r)$ bit cost per dense operation gives
\[
 O\bigl((\sum\delta+kE)r\bigr)
\]
for this reachability work, in addition to the unchanged candidate discovery,
sorting and word summaries. The closure requires $O(r^2)$ bits and reverse
adjacency $O(E\log(r+1))$ bits, apart from container overhead. In a favorable
ordered chain each newly defined vertex has no old ancestors, so propagation
vanishes and the number of graph traversal steps becomes linear. The masks
can still contain $r$ bits; this is not a linear bit-time claim. Neither an
arena-node cap nor a cooperative-work counter alone bounds total process
memory or bit operations.

\paragraph{Recovery-aware adaptive accounting.}
The initial source pass reconstructs, charges and reduces the presentation
before batch planning starts. Its existing conservative accounting includes
block charges for short literal reductions. The old 50,000-work cap can
therefore fail before any search on the 256-crossing circle. The follow-up
does not change these accounting units or the reconstruction algorithm.

For total producer allowance $W$, the host now gives the trial a total
allowance $Q=\lfloor W/4\rfloor$. After successful source reconstruction and
initial root reduction have consumed $R$ work units, the additional search
allowance becomes
\[
 S=\min(50000,Q-R).
\]
The total cap remains in force. A marker records $R$ only after all recovery
work has been charged to the arena. On subsequent local exhaustion the host
deducts the observed trial cost before restarting the established policy
from the original source. If recovery fails before this marker exists,
it conservatively reserves the whole $Q$, or the reported cost if larger.
Aggregate work charges may cross a nominal limit; that reported excess is
also deducted. Both attempts retain one wall deadline, external cancellation
propagates, and mandatory independent replay keeps its existing separate
allowance. Thus the second search receives only the remaining producer work.

The lower-level producer accepts an optional
\texttt{post\_recovery\_work} argument. Its default is \texttt{None}, preserving
the direct caller's full search allowance. The host supplies 50,000. All
public batch flags remain optional and false by default. Increasing the
trial's total cap can spend more work before fallback, so this policy is
not uniformly faster or uniformly stronger under finite limits.

\paragraph{Validation and scope.}
Twelve focused tests pass in 0.693 seconds; the maintained suite passes
1,012 tests with Regina in 215.420 seconds. New tests exercise late ancestor
updates, rejected cycles, changed root order and live sets, both work caps,
strict argument types, and recovery-failure versus search-failure accounting.
Seven hundred test comparisons and 2,250 audit comparisons agree with an
independent literal greedy planner; the audit also compares the pinned old
planner on every state.

The baseline is the complete package at \texttt{47e87a72c23d}. The eighty
source diagrams retain all 240 results in the three earlier modes, and
all eighty direct-batch outcomes, including positive certificate hashes,
agree with the pinned producer.
Every positive audit certificate passes both literal and compressed source
replay. Application, presentation recovery and verifier functions are checked
for identical source, as are the word arena and independent compressed
replay modules. The audit takes 19.573 seconds. Separate 128-, 256- and
512-crossing source trials now finish without fallback at two million total
producer work units; these additional cases are not folded into the common
eighty-diagram coverage claim.

\paragraph{Complete recognition.}
The controlled follow-up uses nineteen fresh source inputs, eight arms
(old and current defaults and batch portfolios, each with an A/A copy),
one warm-up and five shuffled measured rounds. All 760 measured calls and
152 warm-ups complete; the run takes 84.716 seconds. Timing includes the
entire recognition call and mandatory source replay. Ratios are medians
of paired old/new times, not quotients of the marginal time medians.

\begin{center}\small
\begin{tabular}{@{}lrrrrr@{}}
\toprule
Input & default ms & old port. ms & new port. ms & old/new & default/new \\
\midrule
survivor-00 & 17.192 & 15.347 & 15.811 & 0.993 & 1.115 \\
survivor-01 & 17.096 & 46.533 & 41.473 & 0.967 & 0.487 \\
survivor-02 & 15.871 & 10.243 & 10.323 & 0.992 & 1.527 \\
survivor-03 & 11.346 & 43.427 & 43.755 & 0.988 & 0.259 \\
mirror-03 & 24.675 & 6.440 & 6.455 & 1.009 & 3.830 \\
survivor-04 & 10.105 & 31.943 & 31.816 & 1.007 & 0.319 \\
survivor-05 & 19.303 & 16.361 & 16.452 & 0.992 & 1.173 \\
survivor-06 & 7.636 & 5.634 & 5.761 & 0.983 & 1.326 \\
survivor-07 & 19.798 & 21.842 & 21.486 & 1.003 & 0.921 \\
survivor-08 & 7.751 & 7.432 & 7.554 & 0.983 & 1.026 \\
mirror-08 & 20.526 & 16.625 & 16.524 & 1.009 & 1.243 \\
survivor-09 & 7.602 & 5.801 & 5.852 & 0.995 & 1.304 \\
survivor-10 & 6.640 & 5.444 & 5.532 & 0.992 & 1.185 \\
survivor-11 & 7.126 & 5.096 & 5.174 & 0.982 & 1.374 \\
gordian & 1518.937 & 1509.908 & 1623.636 & 0.941 & 0.936 \\
trefoil & 0.052 & 0.053 & 0.053 & 1.004 & 0.995 \\
figure\_eight & 0.060 & 0.060 & 0.059 & 1.003 & 1.013 \\
conway & 0.990 & 1.006 & 0.988 & 1.018 & 1.002 \\
kinoshita\_terasaka & 1.028 & 1.023 & 1.019 & 1.002 & 1.003 \\
\bottomrule
\end{tabular}
\end{center}


The default old/current ratios range from 0.980 to 1.023. Old/new portfolio
ratios range from 0.941 to 1.018, so this sample establishes no broad
whole-input improvement. The four A/A ranges are respectively
0.987--1.016, 0.987--1.031, 0.986--1.016 and 0.967--1.038.
Against the current default, the new optional portfolio still ranges from
0.259 on survivor-03 to 3.830 on mirror-03. Those large gains and regressions
come from the alternative elimination schedule, already present in the
baseline, and should not be attributed to this reachability optimization.

Gordian exposes the price of the larger trial. Recovery costs 39,904 work
units, after which the new search spends its continuation allowance and
falls back. Its trial charge is 89,905 rather than 50,001; total producer
work increases from 1,265,146 to 1,305,050. The fallback proof remains the
same 9,663-byte version-five certificate and its own search still uses
9,025 nodes. Whole-call medians are 1,509.908 ms for the old portfolio and
1,623.636 ms for the new one, with a paired ratio of 0.941 and a new-portfolio
A/A ratio of 1.038. The extra speculative work is real even though noise
affects its timing estimate. Optional status and the existing default are
retained.


\paragraph{Checked group stages on larger circles.}
A separate run bypasses earlier simplification and compares the old and new
direct producers and portfolios, each with an A/A copy. Every successful
call includes source reconstruction, production and mandatory compressed
replay. Seven circle diagrams, one warm-up and five shuffled measured rounds
give 280 measured calls and 56 warm-ups in 45.563 seconds. Of these,
270 measurements and 54 warm-ups complete. All ten incomplete measurements
and both incomplete warm-ups are the old portfolio at 512 crossings, where
fallback hits the unchanged 100,000-node cap. They remain in the evidence;
no old/new completion-time ratio is reported for that portfolio case.

\begin{center}\small
\begin{tabular}{@{}lrrrr@{}}
\toprule
Crossings & direct old/new ms & ratio & portfolio old/new ms & ratio \\
\midrule
8 & 1.409/1.430 & 0.977 & 1.363/1.382 & 0.986 \\
16 & 2.920/2.882 & 1.021 & 2.818/2.824 & 1.004 \\
32 & 6.141/5.797 & 1.060 & 5.915/5.629 & 1.051 \\
64 & 13.458/12.217 & 1.128 & 13.218/12.003 & 1.075 \\
128 & 30.760/24.467 & 1.288 & 377.400/23.901 & 15.998 \\
256 & 74.892/49.422 & 1.515 & 1503.725/50.714 & 30.105 \\
512 & 227.287/101.014 & 2.326 & --/97.486 & -- \\
\bottomrule
\end{tabular}
\end{center}


At 256 crossings the direct checked call improves from 74.892 to 49.422 ms,
with paired ratio 1.515. Producer work drops from 157,823 to 94,320 while
both use 2,544 nodes and the same 12,775-byte certificate. At 512 crossings
the direct ratio is 2.326, work drops from 446,591 to 188,528, and both use
5,104 nodes and the same 25,879-byte certificate. The producer removes
255 or 511 generators in one batch. The reduction is entirely in planning:
these chains require no propagation through old ancestors.

The revised host recovers the direct batch gains previously forfeited by
its trial budget. At 128 crossings its paired ratio against the old host
is 15.998; at 256 it is 30.105. At 256 crossings, recovery consumes 72,677
units and subsequent search 21,643, within the new continuation allowance.
At 512 crossings, recovery consumes 145,125 units and search 43,403;
the new host completes in a median 97.486 ms without fallback. The old
host returns no verdict at its node limit, so this last result demonstrates
a finite-budget coverage difference, not a measured speedup factor.

The stage A/A ranges are 0.979--1.019 for old direct, 0.830--1.023 for
new direct, 0.986--1.031 for completed old portfolios and 0.975--1.030 for
new portfolios. In particular the 128-crossing new-direct control is 0.830,
so its smaller direct timing gain should be treated cautiously. These
controlled stages show where the mechanism helps, while the complete-call
corpus above records the cost and absence of a broad improvement. Neither
set establishes an asymptotic bound for general knot recognition.


Reproduce the evidence using \path{primitive_power_research/elimination_reach.py}
in \path{../fast}, with \texttt{audit}, \texttt{benchmark} and \texttt{stages}
modes. Each run records 140 current source hashes before and after execution,
the pinned baseline, complete certificates, work counters, shuffled orders,
A/A controls and all incomplete outcomes. The optimization retains the
conditional raw-phase theorem above; neither the global depth bound nor
a general quasipolynomial recognizer follows from these measurements.
